% GENERATED FILE. Edit canonical lessons, then run npm run generate:books.
% canonical-source-hash: fnv1a64:1cf03fc78904ae9f
% canonical-lessons: SA-C190-masah, SA-C190-yavah, SA-C190-kharjuram, SA-C190-amalakam, SA-C190-jambuphalam, SA-R190-first-pass-words-from-chapters-182-185, SA-R190-second-pass-words-from-chapters-186-190

\chapter{Black gram, Barley, Date palm's fruit, Amla, Rose apple}
\label{ch:sa-black-gram-barley-date-palm-s-fruit-amla-rose-apple}
\hlchaptermodality{190}

\begin{chapteropening}
\textbf{By the end of this chapter:} \emph{I can say black gram, barley, a date palm's fruit, an amla and a rose apple.}

Complete the last lesson of chapter 190: I can say black gram, barley, a date palm's fruit, an amla and a rose apple.
\end{chapteropening}

\section[\texorpdfstring{\sk{माषः} (māṣaḥ)}{माषः (māṣaḥ)}]{\sk{माषः} (māṣaḥ) — black gram}

\label{lesson:SA-C190-masah}



\begin{quote}
Before the new one: say the Sanskrit for a mat, then the Sanskrit for a carpet.
\end{quote}

\subsection*{You'll want to know: \sk{माषः}}
\textbf{\sk{माषः}} — \emph{māṣaḥ} — \textquotedblleft{}black gram\textquotedblright{}.

\begin{grammarlens}[title={What you've built}]
One more word for this chapter.
\end{grammarlens}

\subsection*{Your turn}
\begin{itemize}
\raggedright
  \item \emph{Say it:} \emph{māṣaḥ}
  \item \emph{Say it:} \emph{māṣaḥ}, once more
  \item \emph{Say it:} say \emph{māṣaḥ}
  \item \emph{Recall:} say the Sanskrit for a mat, then the Sanskrit for a carpet, then say \emph{māṣaḥ} again
\end{itemize}

\begin{tcolorbox}[breakable,colback=teal!4,colframe=teal!35!black,title={Before you move on}]
What does \sk{माषः} mean? (\textquotedblleft{}Black gram\textquotedblright{}.) Say it once more.
\end{tcolorbox}
\section[\texorpdfstring{\sk{यवः} (yavaḥ)}{यवः (yavaḥ)}]{\sk{यवः} (yavaḥ) — barley}

\label{lesson:SA-C190-yavah}



\begin{quote}
Before the new one: say the Sanskrit for a carpet, then the Sanskrit for black gram.
\end{quote}

\subsection*{You'll want to know: \sk{यवः}}
\textbf{\sk{यवः}} — \emph{yavaḥ} — \textquotedblleft{}barley\textquotedblright{}.

\begin{grammarlens}[title={What you've built}]
One more word for this chapter.
\end{grammarlens}

\subsection*{Your turn}
\begin{itemize}
\raggedright
  \item \emph{Say it:} \emph{yavaḥ}
  \item \emph{Say it:} \emph{yavaḥ}, once more
  \item \emph{Say it:} say \emph{yavaḥ}
  \item \emph{Recall:} say the Sanskrit for a carpet, then the Sanskrit for black gram, then say \emph{yavaḥ} again
\end{itemize}

\begin{tcolorbox}[breakable,colback=teal!4,colframe=teal!35!black,title={Before you move on}]
What does \sk{यवः} mean? (\textquotedblleft{}Barley\textquotedblright{}.) Say it once more.
\end{tcolorbox}
\section[\texorpdfstring{\sk{खर्जूरम्} (kharjūram)}{खर्जूरम् (kharjūram)}]{\sk{खर्जूरम्} (kharjūram) — a date palm's fruit}

\label{lesson:SA-C190-kharjuram}



\begin{quote}
Before the new one: say the Sanskrit for black gram, then the Sanskrit for barley.
\end{quote}

\subsection*{You'll want to know: \sk{खर्जूरम्}}
\textbf{\sk{खर्जूरम्}} — \emph{kharjūram} — \textquotedblleft{}a date palm's fruit\textquotedblright{}.

\begin{grammarlens}[title={What you've built}]
One more word for this chapter.
\end{grammarlens}

\subsection*{Your turn}
\begin{itemize}
\raggedright
  \item \emph{Say it:} \emph{kharjūram}
  \item \emph{Say it:} \emph{kharjūram}, once more
  \item \emph{Say it:} say \emph{kharjūram}
  \item \emph{Recall:} say the Sanskrit for black gram, then the Sanskrit for barley, then say \emph{kharjūram} again
\end{itemize}

\begin{tcolorbox}[breakable,colback=teal!4,colframe=teal!35!black,title={Before you move on}]
What does \sk{खर्जूरम्} mean? (\textquotedblleft{}A date palm's fruit\textquotedblright{}.) Say it once more.
\end{tcolorbox}
\section[\texorpdfstring{\sk{आमलकम्} (āmalakam)}{आमलकम् (āmalakam)}]{\sk{आमलकम्} (āmalakam) — an amla, gooseberry}

\label{lesson:SA-C190-amalakam}



\begin{quote}
Before the new one: say the Sanskrit for barley, then the Sanskrit for a date palm's fruit.
\end{quote}

\subsection*{You'll want to know: \sk{आमलकम्}}
\textbf{\sk{आमलकम्}} — \emph{āmalakam} — \textquotedblleft{}an amla, gooseberry\textquotedblright{}.

\begin{grammarlens}[title={What you've built}]
One more word for this chapter.
\end{grammarlens}

\subsection*{Your turn}
\begin{itemize}
\raggedright
  \item \emph{Say it:} \emph{āmalakam}
  \item \emph{Say it:} \emph{āmalakam}, once more
  \item \emph{Say it:} say \emph{āmalakam}
  \item \emph{Recall:} say the Sanskrit for barley, then the Sanskrit for a date palm's fruit, then say \emph{āmalakam} again
\end{itemize}

\begin{tcolorbox}[breakable,colback=teal!4,colframe=teal!35!black,title={Before you move on}]
What does \sk{आमलकम्} mean? (\textquotedblleft{}An amla, gooseberry\textquotedblright{}.) Say it once more.
\end{tcolorbox}
\section[\texorpdfstring{\sk{जम्बूफलम्} (jambūphalam)}{जम्बूफलम् (jambūphalam)}]{\sk{जम्बूफलम्} (jambūphalam) — a rose apple}

\label{lesson:SA-C190-jambuphalam}



\begin{quote}
Before the new one: say the Sanskrit for a date palm's fruit, then the Sanskrit for an amla.
\end{quote}

\subsection*{You'll want to know: \sk{जम्बूफलम्}}
\textbf{\sk{जम्बूफलम्}} — \emph{jambūphalam} — \textquotedblleft{}a rose apple\textquotedblright{}.

\begin{grammarlens}[title={What you've built}]
One more word for this chapter.
\end{grammarlens}

\subsection*{Your turn}
\begin{itemize}
\raggedright
  \item \emph{Say it:} \emph{jambūphalam}
  \item \emph{Say it:} \emph{jambūphalam}, once more
  \item \emph{Say it:} say \emph{jambūphalam}
  \item \emph{Recall:} say the Sanskrit for a date palm's fruit, then the Sanskrit for an amla, then say \emph{jambūphalam} again
\end{itemize}

\begin{tcolorbox}[breakable,colback=teal!4,colframe=teal!35!black,title={Before you move on}]
What does \sk{जम्बूफलम्} mean? (\textquotedblleft{}A rose apple\textquotedblright{}.) Say it once more.
\end{tcolorbox}
\section[(dialogue)]{Words from chapters 182-185 — a first pass}

\label{lesson:SA-R190-first-pass-words-from-chapters-182-185}



\begin{quote}
Say the words for an amla and a rose apple.
\end{quote}

\subsection*{Your turn}
Say these aloud:

\begin{itemize}
\raggedright
  \item wonder, grief, contentment, gratitude, memory
  \item an age, an auspicious hour, an hour, for a long time, after a long time
  \item from afar, close by, between, loudly, softly
  \item how much?, of what kind?, for, surely, certainly
\end{itemize}

\begin{tcolorbox}[breakable,colback=teal!4,colframe=teal!35!black,title={Before you move on}]
Wonder? (\textbf{\sk{विस्मयः}}, \emph{vismayaḥ}.) Grief? (\textbf{\sk{शोकः}}, \emph{śokaḥ}.) Contentment? (\textbf{\sk{सन्तोषः}}, \emph{santoṣaḥ}.) Gratitude? (\textbf{\sk{कृतज्ञता}}, \emph{kṛtajñatā}.) Memory? (\textbf{\sk{स्मृतिः}}, \emph{smṛtiḥ}.) An age? (\textbf{\sk{युगम्}}, \emph{yugam}.) An auspicious hour? (\textbf{\sk{मुहूर्तः}}, \emph{muhūrtaḥ}.) An hour? (\textbf{\sk{होरा}}, \emph{horā}.) For a long time? (\textbf{\sk{चिरम्}}, \emph{ciram}.) After a long time? (\textbf{\sk{चिरात्}}, \emph{cirāt}.) From afar? (\textbf{\sk{दूरतः}}, \emph{dūrataḥ}.) Close by? (\textbf{\sk{निकटे}}, \emph{nikaṭe}.) Between? (\textbf{\sk{अन्तरा}}, \emph{antarā}.) Loudly? (\textbf{\sk{उच्चैः}}, \emph{uccaiḥ}.) Softly? (\textbf{\sk{नीचैः}}, \emph{nīcaiḥ}.) How much?? (\textbf{\sk{कियत्}}, \emph{kiyat}.) Of what kind?? (\textbf{\sk{कीदृशः}}, \emph{kīdṛśaḥ}.) For? (\textbf{\sk{हि}}, \emph{hi}.) Surely? (\textbf{\sk{नूनम्}}, \emph{nūnam}.) Certainly? (\textbf{\sk{अवश्यम्}}, \emph{avaśyam}.)
\end{tcolorbox}
\section[(dialogue)]{Words from chapters 186-190 — a second pass}

\label{lesson:SA-R190-second-pass-words-from-chapters-186-190}



\begin{quote}
Say the words for tender and dwarfish.
\end{quote}

\subsection*{Your turn}
Say these aloud:

\begin{itemize}
\raggedright
  \item tender, dwarfish, vast, fine, steady
  \item kind, generous, stingy, truthful, adept
  \item thirsty, ready, busy, spotless, bright
  \item a trunk, a hanging basket, a basket, a mat, a carpet
  \item black gram, barley, a date palm's fruit, an amla, a rose apple
\end{itemize}

\begin{tcolorbox}[breakable,colback=teal!4,colframe=teal!35!black,title={Before you move on}]
Tender? (\textbf{\sk{कोमलः}}, \emph{komalaḥ}.) Dwarfish? (\textbf{\sk{वामनः}}, \emph{vāmanaḥ}.) Vast? (\textbf{\sk{विशालः}}, \emph{viśālaḥ}.) Fine? (\textbf{\sk{सूक्ष्मः}}, \emph{sūkṣmaḥ}.) Steady? (\textbf{\sk{स्थिरः}}, \emph{sthiraḥ}.) Kind? (\textbf{\sk{दयालुः}}, \emph{dayāluḥ}.) Generous? (\textbf{\sk{उदारः}}, \emph{udāraḥ}.) Stingy? (\textbf{\sk{कृपणः}}, \emph{kṛpaṇaḥ}.) Truthful? (\textbf{\sk{सत्यवादी}}, \emph{satyavādī}.) Adept? (\textbf{\sk{पटुः}}, \emph{paṭuḥ}.) Thirsty? (\textbf{\sk{तृषितः}}, \emph{tṛṣitaḥ}.) Ready? (\textbf{\sk{सज्जः}}, \emph{sajjaḥ}.) Busy? (\textbf{\sk{व्यस्तः}}, \emph{vyastaḥ}.) Spotless? (\textbf{\sk{निर्मलः}}, \emph{nirmalaḥ}.) Bright? (\textbf{\sk{उज्ज्वलः}}, \emph{ujjvalaḥ}.) A trunk? (\textbf{\sk{मञ्जूषा}}, \emph{mañjūṣā}.) A hanging basket? (\textbf{\sk{शिक्यम्}}, \emph{śikyam}.) A basket? (\textbf{\sk{पिटकम्}}, \emph{piṭakam}.) A mat? (\textbf{\sk{कटः}}, \emph{kaṭaḥ}.) A carpet? (\textbf{\sk{आस्तरणम्}}, \emph{āstaraṇam}.) Black gram? (\textbf{\sk{माषः}}, \emph{māṣaḥ}.) Barley? (\textbf{\sk{यवः}}, \emph{yavaḥ}.) A date palm's fruit? (\textbf{\sk{खर्जूरम्}}, \emph{kharjūram}.) An amla? (\textbf{\sk{आमलकम्}}, \emph{āmalakam}.) A rose apple? (\textbf{\sk{जम्बूफलम्}}, \emph{jambūphalam}.)
\end{tcolorbox}
